\documentclass[11pt]{article}
\usepackage[a4paper, top=18mm, bottom=18mm, left=20mm, right=20mm]{geometry}
\usepackage[T2A]{fontenc}
\usepackage[utf8]{inputenc}
\usepackage[ukrainian]{babel}
\usepackage{graphicx}
\usepackage[export]{adjustbox}
\usepackage{amsmath}
\usepackage{amssymb}
\graphicspath{{/app/Quiz/Scripts/2023/ProgAll/15BinTree/}{/app/Quiz/Scripts/2020/Mod2020_4/BinTree/}}
\setlength{\parindent}{0pt}
\setlength{\parskip}{4pt}
\pagestyle{empty}
\begin{document}
%begin
{\large\textbf{Контрольна робота}}

\textbf{Варіант \No\,1}\hfill \textit{ПІБ:}\,\underline{\hspace{60mm}}
\vspace{4mm}

%beginitem
1. Що буде виведено за виконання фрагменту коду?

%code
1**1.0*False
%code

%enditem

%beginitem
2. Що буде виведено за виконання фрагменту коду?
%code

x,y,z=5,4,0
z,y,x,y=x,y,8,9
print(x,y,z)
%code

%enditem

%beginitem
3. 
try:
    print(5, end=';')
    print(int('c4'), end=';')
    print(0, end=';')
except ZeroDivisionError:
    print(1, end=';')
except KeyboardInterrupt:
    print(6, end=';')
except:
    print('ERR', end=';')
else:
    print(9, end=';')
finally:
    print(3)

%enditem

%beginitem
4. %goodpath:Scripts/2018/Mod2018_1/01-good.txt
Відмітити все, що є однією правильною лексемою:
( 1 ) x*y

( 2 ) 83x

( 3 ) +

( 4 ) for

%enditem

%beginitem
5. 
1**1.0*False
%enditem

%beginitem
6. Що буде виведено за виконання фрагменту коду?

%code
t=5
while t<8:
    t+=1
print(t, end=' ')
%code


%enditem

%beginitem
7. Що буде виведено за виконання фрагменту коду?

%code
print('R', end=' ')
f=[10,11,12,13,14]
f.insert(-4, '12')
print(f)
%code


%enditem

%beginitem
8. Що буде виведено за виконання фрагменту коду?

%code
cout << 6 * 6 * 6 * 6;
%code

%enditem

%beginitem
9. Що буде виведено за виконання фрагменту коду?

%code
try:
    t = 8
    while t>5:
        t += -3
    print(t, end=';')
except:
    print('Z')

%code

%enditem

%beginitem
10. 
try:
    res = "True"!="4.0"
    if isinstance(res, bool):
        print(f'bool;{res}')
    elif isinstance(res, int):
        print(f'int;{res}')
    elif isinstance(res, float):
        print(f'float;{res:.2f}')
    else:
        print('???')
except:
    print('Z')

%enditem

%beginitem
11. Що буде виведено за виконання фрагменту коду?

%code
print('R', end=' ')
f = ('a','b','c','d','e',0,1,2,4)
print(f[12])
%code


%enditem

%beginitem
12. Що буде виведено за виконання фрагменту коду?
%code
if (5 != 5)
    cout << "y";
else
    cout << "w";
%code


%enditem

%beginitem
13. Що буде виведено за виконання фрагменту коду?

%code
cout << (6 * 6 * 6 * 6);
%code

%enditem

%beginitem
14. 
Записати ВИРАЗ, значенням якого є множина, що складається з квадратівцілих чисел від 13 до 2003 (включно), що не діляться на 5.
%enditem

%beginitem
15. Що буде виведено за виконання фрагменту коду?

%code
print(notTrue!=7 and not6==7.0 and 6.0<2)
%code

%enditem

%beginitem
16. Що буде виведено за виконання фрагменту коду?

%code
print('R', end=' ')
f = ('0','1','2','3','4','5','6','7','8','9')
print(f[:10:2])
%code


%enditem

%beginitem
17. 
try:
    for f in range(10, 10, 2):
        if f>8:
            continue
        print(f, end=';')
        if f>8:
            break
    else:
        print(f,end=';')
    print(f)
except:
    print('Z')

%enditem

%beginitem
18. Що буде виведено за виконання фрагменту коду?
%Оригинал был утерян. Требует отладки!
%code
_c = 0

class C:
    _c = 1
    
    def __init__(self):
        self._c = 2
        _c = 3
        C._c = 4

obj = C()
obj._c = 5
try:
    print(_c, end=' ')
    print(obj._c, end=' ')
    print(C._c, end=' ')
    print(obj._c, end=' ')
    print(obj._C__c)
except:
    print('ERR')
%code

%enditem

%beginitem
19. %goodpath:Scripts/2019/Mod2019_1/Lex/03-good.txt
Відмітити все, що є коректним оператором С++:
( 1 ) (2+2j)//2

( 2 ) input()*1.2

( 3 ) 3.14%2

( 4 ) x=1%3

%enditem

%beginitem
20. 

print('R', end=' ')
f=('0','1','2','3','4','5','6','7','8','9')
print(f[:10:2])


%enditem

%beginitem
21. 
try:
    res = "True"!="4.0"
    if isinstance(res, bool):
        print(f'bool;{res}')
    elif isinstance(res, int):
        print(f'int;{res}')
    elif isinstance(res, float):
        print(f'float;{res:.2f}')
    else:
        print('???')
except:
    print('ERR')

%enditem

%beginitem
22. Що буде виведено за виконання фрагменту коду?

%code
try:
    t = 5
    while t<8:
        t += 1
    print(t, end=';')
except:
    print('Z')

%code

%enditem

%beginitem
23. 
Змінні \kw{next} та \kw{prev} вузла списку позначають наступний та попередній за ним вузол відповідно. Починаючи з голови, у список записано послідовні цілі числа від~$-19$ до~$39$. Нехай змінна \kw{p} позначає вузол, що зберігає значення~$-6$.
Яке число зберігається у вузлі \kw{p.prev.prev.prev.prev.prev} ?

%enditem

%beginitem
24. 
a = 5
b = 4
c = 0

def f(a):
global c    a = 4
    b = 5
    c = 3
    return a + b + c

a = 1
b = 6
c = 9

try:
    print(f(b), a, b, c, sep=';')
except:
    print('Z', a, b, c, sep=';')

%enditem

%beginitem
25. Що буде виведено за виконання фрагменту коду?

%code
t=8
while t>5:
    t+=-3
print(t, end=' ')
%code


%enditem

%beginitem
26. 
В умовах попередньої задачі було виконано p->prev=p->prev->prev;

Чому дорівнює значення p->prev->prev->prev->prev->n ?
%enditem

%beginitem
27. Що буде виведено за виконання фрагменту коду?

%code
cout << 6 * 6 * 6 * 6;
%code

%enditem

%beginitem
28. Що буде виведено за виконання фрагменту коду?

%code
f=[10,11,12,13]
f.insert(-3, '12')
print(f)
%code


%enditem

%beginitem
29. 
Значенням змінної \kw{a} є цілочисловий масив типу $numpy.ndarray$ розмірів $6{\times}6$. На скільки збільшиться сума елементів масиву \kw{a} після виконання
\begin{verbatim}
a[2] += 1
\end{verbatim}


Якщо код некоректний, то в якості відповіді зазначити ERROR.



%enditem

%beginitem
30. Що буде виведено за виконання фрагменту коду?
%code
if (5 != 5)
    cout << "y";
else
    cout << "y";
%code


%enditem

%beginitem
31. 

def f(n):
    print("f", end="");
    return n

print(f(-6) and f(7))

%enditem

%beginitem
32. 
try:
    print(5, end=';')
    print(int('c4'), end=';')
    print(0, end=';')
except ZeroDivisionError:
    print(1, end=';')
except KeyboardInterrupt:
    print(6, end=';')
except:
    print('Z', end=';')
else:
    print(9, end=';')
finally:
    print(3)

%enditem

%beginitem
33. Що буде виведено за виконання фрагменту коду?

%code
print('R', end=' ')
f=('0','1','2','3','4','5','6','7','8','9')
print(f[:10:2])
%code


%enditem

%beginitem
34. Що буде виведено за виконання фрагменту коду?
code
c = 0
n = 1

class C:
    c = 2
    
    def __init__(self):
global c        self.c = 3
        c = 4
        C.c = 5

obj = C()
print(c, n, C.c, obj.c, sep=' ')
code

%enditem

%beginitem
35. 
Записати ВИРАЗ, значенням якого є список, що складається з цілих чисел від 13 до 2003 (включно), причому спочатку за зростанням записано всі, що не діляться на 5, а потім решту за зростанням.
%enditem

%beginitem
36. 
notTrue!=7 and not6==7.0 and 6.0<2
%enditem

%beginitem
37. 
a = 5
b = 4
c = 0

def f(a):
global c    a = 4
    b = 5
    c = 3
    return a + b + c

a = 1
b = 6
c = 9

try:
    print(f(b), a, b, c, sep=';')
except:
    print('ERR', a, b, c, sep=';')

%enditem

%beginitem
38. Що буде виведено за виконання фрагменту коду?

%code
for f in range(-7,-7+5,2):
    if f>8:
        break
    print(f, end=' ')
    f=-7
    if f>8:
        break
else:
    print(f, end=' ')
print(f, end=' ')
%code

%enditem

%beginitem
39. 

\begin{center}
	\adjustimage{max size={0.8\linewidth}{0.8\paperheight}}
	{../15BinTree/trees_exp/121.png}
\end{center}
Указати послідовність міток вершин дерева, зображеного на рис., що утвориться в результаті обходу дерева в прямому порядку. Мітки записувати через пробіл.\par Алгоритм починає роботу з кореня (на рис. це верхня вершина).
%answer:121 прямому

%enditem

%beginitem
40. Що буде виведено за виконання фрагменту коду?

%code
t=8
while t>5:
    t+=-3
print(t)
%code


%enditem

%beginitem
41. 

x,y,z=5,4,0
z,y,x,y=x,y,8,9
print(x,y,z)

%enditem

%beginitem
42. 
6/6//6*6
%enditem

%beginitem
43. Що буде виведено за виконання фрагменту коду?
%code
a, b, c = 8, 9, 6
def f(a, b=7, c):
    print(a, b, c, end=" ")

f(5, c=4, b=0)
print(a, b, c)
%code

%enditem

%beginitem
44. 
Записати ВИРАЗ, значенням якого є словник, ключами якого є цілі числа від 13 до 2003 (включно), що не діляться на 5, а ключам відповідають їх квадрати 
%enditem

%beginitem
45. Що буде виведено за виконання фрагменту коду?

%code
print(6/6//6*6)
%code

%enditem

%beginitem
46. 
def f(arg, n):
    n = 5
    tmp = arg
    tmp[-4:-4] = 'ac'
    return len(tmp)>3

try:
    f = [10,11,12,13,14]
    n = 4
    f(f, n)
    print(n, end=';')
    for el in f:
        print(el, end=';')
except:
    print('Z')

%enditem

%beginitem
47. Що буде виведено за виконання фрагменту коду?

%code
notTrue!=7 and not6==7.0 and 6.0<2
%code

%enditem

%beginitem
48. Що буде виведено за виконання фрагменту коду?

%code
#include <iostream>
using namespace std;

int f(int a, int b){
    int c = 61;
    if (b != -2)
        return 5;
    if (b == 5)
         c = 4;
    else 
        c = 0;
    return c;
}

int main(){
    cout << f(-6, -6);
    return 0;
}
%code

%enditem

%beginitem
49. Що буде виведено за виконання фрагменту коду?

%code
6/6//6*6
%code

%enditem

%beginitem
50. Що буде виведено за виконання фрагменту коду?

%code
a,b,c=5,4,0
def f(a):
    a=4
    b+=5
    c=3
    return a+b+c

a,b,c=1,6,9
print(f(b),a,b,c)
%code

%enditem

%beginitem
51. 
def f(n):
    print('f', end=';')
    return n

try:
    print(f(-6) and f(7))
except:
    print('Z')

%enditem

%beginitem
52. 
try:
    res = notTrue!=7 and not6==7.0 and 6.0<2
    if isinstance(res, bool):
        print(f'bool;{res}')
    elif isinstance(res, int):
        print(f'int;{res}')
    elif isinstance(res, float):
        print(f'float;{res:.2f}')
    else:
        print('???')
except:
    print('ERR')

%enditem

%beginitem
53. 
Рядки журналу через двокрапку містять ім'я користувача (ідентифікатор), час події,тип події, код події, наприклад
\begin{verbatim}
s="username : 2022-05-05 13:26 : ERROR : 404"
\end{verbatim}
у коді 
\begin{verbatim}
res = re.fullmatch('???', s) 
s = ''
code_str = ??? 
\end{verbatim}
чим треба замінити кожен з ???, щоб значенням code\_str був код події? 



%enditem

%beginitem
54. %goodpath:Scripts/2018/Mod2018_1/03-good.txt
Відмітити все, що є коректним оператором С++:
( 1 ) 'a'=3;

( 2 ) int f(int a, int 4);

( 3 ) bool r=('N'>'5');

( 4 ) int f(int &n, int m);

%enditem

%beginitem
55. Що буде виведено за виконання фрагменту коду?

%code
#include <iostream>
using namespace std;

bool f(int n){
    cout<<"f";
    return n;
}

int main(){
    cout<<(f(-6) && f(7));
    return 0;
}
%code

%enditem

%beginitem
56. 

\begin{center}
	\adjustimage{max size={0.8\linewidth}{0.8\paperheight}}
	{../BinTree/trees_exp/121.png}
\end{center}
Указати послідовність міток вершин дерева, зображеного на рис., що утвориться в результаті обходу дерева в прямому порядку. Мітки записувати через пробіл.\par Обхід дерева починається з кореня (на рис. це верхня вершина).
%answer:121 прямому

%enditem

%beginitem
57. 
def f(n):
    print('f', end='')
    return n

try:
    print(f(-6) and f(7))
except:
    print('ERR')

%enditem

%beginitem
58. Що буде виведено за виконання фрагменту коду?

%code
"True"!="4.0"
%code

%enditem

%beginitem
59. Що буде виведено за виконання фрагменту коду?

%code
#include <iostream>

int f(int a, int b){
    int c = 61;
    if (b != -2)
        return 5;
    if (b == 5)
         c = 4;
    else 
        c = 0;
    return c;
}

int main(){
    std::cout << f(-6, -6);
    return 0;
}
%code

%enditem

%beginitem
60. 
Кожен рядок текстового файлу містить по два дійсних числа, розділені білими символами.
Написати функцію, що для заданого текстового файлу будує новий текстовий файл, рядки якого крім зображень дйсних чисел ще містять результат обчислення на них функції $f$. У вихідному файлі зображення чисел мають розділятися пробілами. 
$$f(x, y) =\frac{\sqrt x + 34}{(\ctg x + 0.4) (y + 60)}$$ 
Якість коду також оцінюється.


%enditem

%beginitem
61. Що буде виведено за виконання фрагменту коду?

%code
print("True" != "4.0")
%code

%enditem

%beginitem
62. Що буде виведено за виконання фрагменту коду?

%code
#include <iostream>

bool f(int n){
    std::cout<<"f";
    return n;
}

int main(){
    std::cout<<(f(-6) && f(7));
    return 0;
}
%code

%enditem

%beginitem
63. 

print('R', end=' ')
f=[10,11,12,13,14]
f[-4:-4]='11'
print(f)


%enditem

%beginitem
64. Що буде виведено за виконання фрагменту коду?

%code
f = ('a','b','c','d','e',0,1,2,4)
print(f[12])
%code


%enditem

%beginitem
65. 

print('R', end=' ')
for f in range(10,10,2):
    if f>8:
        continue
        print(f, end=' ')
    if f>8:
        break
    else:
        print(f, end=' ')
    print(f)

%enditem

%beginitem
66. Що буде виведено за виконання фрагменту коду?
%code
c = 0
n = 1

class C:
    c = 2
    
    def __init__(self):
global c        self.c = 3
        c = 4
        C.c = 5

obj = C()
try:
    print(c, end=' ')
    print(n, end=' ')
    print(C.c, end=' ')
    print(obj.c)
except:
    print('ERR')
%code

%enditem

%beginitem
67. 
try:
    res = notTrue!=7 and not6==7.0 and 6.0<2
    if isinstance(res, bool):
        print(f'bool;{res}')
    elif isinstance(res, int):
        print(f'int;{res}')
    elif isinstance(res, float):
        print(f'float;{res:.2f}')
    else:
        print('???')
except:
    print('Z')

%enditem

%beginitem
68. Що буде виведено за виконання фрагменту коду?

%code
print('R', end=' ')
for f in range(-7,-7,2):
    if f>8:
        continue
        print(f, end=' ')
    if f>8:
        break
    else:
        print(f, end=' ')
    print(f, end=' ')
%code

%enditem

%beginitem
69. Що буде виведено за виконання фрагменту коду?
%code
a,b,c=8,9,6
def f(a,b=7,c):
    print(a,b,c,end=" ")

f(5,c=4,b=0)
print(a,b,c)
%code

%enditem

%beginitem
70. %goodpath:Scripts/2019/Mod2019_1/Lex/01-good.txt
Відмітити все, що є однією правильною лексемою:
( 1 ) [1]

( 2 ) 2int

( 3 ) False

( 4 ) *

%enditem

%beginitem
71. Що буде виведено за виконання фрагменту коду?

%code
def f(a,b):
    c=61
    if b!=-2:
        return 5
    if b==5:
         c=4
    else: 
        c=0
    return c

print(f(-6,-6))
%code

%enditem

%beginitem
72. Що буде виведено за виконання фрагменту коду?
%code
#include <iostream>
using namespace std;

int f(int &x, int &y){
    x = 4;
    y+= 6;
    return x;
}

int main(){
    int a = 7, b = 9;
    a = f(a, a);
    cout << a << ":" << b <<':';
    {
        int a = 2, b = 8;
        cout << ((a>6) && ((b+=1) > 6))
             << a << ":" << b << ":";
    }
    {
        inta = 3;
        cout << a << ':';
    }
    cout << a << ":" << b << endl;
    return 0;
}
%code


%enditem

%beginitem
73. Що буде виведено за виконання фрагменту коду?
%code
#include <iostream>
using namespace std;

int f(int &x, int &y){
    x = 4;
    y+= 6;
    return x;
}

int main(){
    int a = 7, b = 9;
    a = f(a, a);
    cout << a << ":" << b <<':';
    {
        int a = 2, b = 8;
        cout << ((a>6) && ((b+=1) > 6)) << ':';
        cout << a << ":" << b << ':';
    }
    {
        inta = 3;
        cout << a << ':';
    }
    cout << a << ":" << b << endl;
    return 0;
}
%code


%enditem

%beginitem
74. 
$a$ є цілочисловим масивом типу $numpy.ndarray$ розмірів $4{\times}4$. Сума елементів масиву $a$ дорівнює 45. Чому дорівнює сума елементів масиву $a$ після виконання 
\begin{verbatim}
a += 45
\end{verbatim}


Якщо код некоректний, то в якості відповіді зазначити ERROR.



%enditem

%beginitem
75. 
Задано тип вузла списку
struct Node \{int n; Node *prev, *next;\};
У списку зберігаються послідовні цілі числа від -19 до 39. Вказівник p встановлено на вузол, в якому зберігається число -6.
Чому дорівнює значення p->prev->prev->prev->prev->prev->n ?
%enditem

%beginitem
76. 

def f(a,b):
    c=61
    if b!=-2:
        return 5
    if b==5:
         c=4
    else: 
        c=0
    return c

print(f(-6,-6))
%enditem

%beginitem
77. Що буде виведено за виконання фрагменту коду?

%code
cout << (6 * 6 * 6 * 6);
%code

%enditem

%beginitem
78. 

a,b,c=5,4,0
def f(a):
global c    a=4
    b+=5
    c=3
    return a+b+c

a,b,c=1,6,9
print(f(b),a,b,c)

%enditem

%beginitem
79. Що буде виведено за виконання фрагменту коду?
%code
for f in range(10,10,2):
    if f>8:
        continue
        print(f,end=' ')
    if f>8:
        break
else:
    print(f, end=' ')
print(f, end=' ')
%code

%enditem

%beginitem
80. 
Написати програму, що запитує у користувача значення дійсних параметрів $x$ та $y$, обчислює для них значення виразу 
$$\frac{\sqrt x + 34}{(\ctg x + 0.4) (y + 60)}$$ 
та повiдомляє введенi данi та результат обчислення.


%enditem

%beginitem
81. 
try:
    res = 6/6//6*6
    if isinstance(res, bool):
        print(f'bool;{res}')
    elif isinstance(res, int):
        print(f'int;{res}')
    elif isinstance(res, float):
        print(f'float;{res:.2f}')
    else:
        print('???')
except:
    print('ERR')

%enditem

%beginitem
82. Що буде виведено за виконання фрагменту коду?

%code
#include <iostream>
using namespace std;

int f(int c){
    int z = 61;
    if (c != -2) 
        return 5;
    if (c == 5)
         z = 4;
    else
         z = 0;
    return z;
}

int main(){
    cout << f(-6);
    return 0;
}
%code


%enditem

%beginitem
83. Що буде виведено за виконання фрагменту коду?

%code
print('R', end=' ')
t=8
while t>5:
    t+=-3
print(t)
%code


%enditem

%beginitem
84. 
try:
    t = 5
    while t<8:
        t += 1
    print(t, end=';')
except:
    print('Z')

%enditem

%beginitem
85. 

a,b,c=8,9,6
def f(a,b=7,c):
    print(a,b,c,end="")

f(5,c=4,b=0)
print(a,b,c)


%enditem

%beginitem
86. Що буде виведено за виконання фрагменту коду?

%code
print('R', end=' ')
t=5
while t<8:
    t+=1
print(t)
%code


%enditem

%beginitem
87. Що буде виведено за виконання фрагменту коду?

%code
#include <iostream>
using namespace std;

int a = 5, b = 4, c = 0;

int f(int &a){
int c;    a = 4;
    b = 5;
    c = 3;
    return a + b + c;
}

int main(){
    inta = 1;
    intb = 6;
    intc = 9;
    cout << f(b) << ':';
    cout << a << ':' << b << ':' << c;
    return 0; 
}
%code

%enditem

%beginitem
88. 

def f(c):
    z=61
    if c!=-2: 
        return 5
    if c==5:
         z=4
    else:
         z=0
    return z

print(f(-6))

%enditem

%beginitem
89. 

print('R', end=' ')
f=[10,11,12,13,14]
f.insert(-4, '12')
print(f)


%enditem

%beginitem
90. 

print('R', end=' ')
t=8
while t>5:
    t+=-3
print(t)


%enditem

%beginitem
91. 
f=[10,11,12,13,14]; f[-4:-4]='11'; print(f)
%enditem

%beginitem
92. 

print('R', end=' ')
t=5
while t<8:
    t+=1
print(t)


%enditem

%beginitem
93. Що буде виведено за виконання фрагменту коду?

%code
#include <iostream>

int f(int c){
    int z = 61;
    if (c != -2) 
        return 5;
    if (c == 5)
         z = 4;
    else
         z = 0;
    return z;
}

int main(){
    std::cout << f(-6);
    return 0;
}
%code


%enditem

%beginitem
94. 
Змінні $next$ та $prev$ вузла списку позначають наступний та попередній за ним вузол відповідно. Починаючи з голови, у список записано послідовні цілі числа від $-19$ до $39$. Нехай змінна $p$ позначає вузол, що зберігає значення $-6$.
Яке значення зберігається у вузлі $p.prev.prev.prev.prev.prev$ ?

%enditem

%beginitem
95. 
try:
    res = 1**1.0*False
    if isinstance(res, bool):
        print(f'bool;{res}')
    elif isinstance(res, int):
        print(f'int;{res}')
    elif isinstance(res, float):
        print(f'float;{res:.2f}')
    else:
        print('???')
except:
    print('Z')

%enditem

%beginitem
96. Що буде виведено за виконання фрагменту коду?

%code
print(True >= 7 != 6 is 7.0)
%code

%enditem

%beginitem
97. 
Написати програму, що запитує у користувача значення дійсних параметрів $x$ та $y$, обчислює для них значення виразу 
$$\frac{\sqrt x + 34}{(\ctg x + 0.4) (y + 60)}$$ 
та повiдомляє введенi данi та результат обчислення.
 Оцінюється правильність та якість коду.


%enditem

%beginitem
98. Що буде виведено за виконання фрагменту коду?

%code
cout << (!true!=7 && !6== 7.0 && 6.0<2);
%code

%enditem

%beginitem
99. Що буде виведено за виконання фрагменту коду?
%code

def f(n):
    print("f", end="");
    return n

print(f(-6) and f(7))
%code

%enditem

%beginitem
100. Що буде виведено за виконання фрагменту коду?

%code
a = 5
b = 4
c = 0

def f(a):
global c    a=4
    b+=5
    c=3
    return a+b+c

a, b, c = 1, 6, 9
try:
    print(f(b), a, b, c, sep=';')
except:
    print('Z', a, b, c, sep=';')
%code

%enditem

%end
\newpage
\end{document}

